\section{A corpus nobody's product owns}
\label{sec:28.5}
The book has said what a public culture would do: make work free to read and learn from, by humans and by machines, and pay the people who make it by something other than the policing of copies. Its note said a polity can pay them by appropriation. Here is the institution. An endowment pays makers for work that enters a public corpus, held by the Library of Congress, free to read and learn from. The maker is named with the work, and paid for putting it there. The formation authority draws on the corpus for the material a model is first formed on, so that what forms the public's models isn't whatever a product happened to collect.

The United States has written the purpose down before. The National Foundation on the Arts and the Humanities Act of 1965 finds that “an advanced civilization must not limit its efforts to science and technology alone,” and that “democracy demands wisdom and vision in its citizens” and “must therefore foster and support a form of education, and access to the arts and the humanities, designed to make people of all backgrounds and wherever located masters of their technology and not its unthinking servants” \autocite{us1965artsact}.

The same Act shows how a body like this is captured. In 1990 Congress told the arts endowment to judge applications “taking into consideration general standards of decency and respect for the diverse beliefs and values of the American public,” and in 1998 the Supreme Court upheld the clause on its face, reading it as a consideration and not a rule \autocite{us1965artsact,scotus1998finley}. A clause like that is a handle an administration can turn without amending anything. So the endowment carries no content criteria. The corpus leaves out what any corpus leaves out, instructions for weapons and the sexual abuse of children, and nothing else on its content.

Payment is the guarantee's (Appendix~\ref{sec:C}): a panel makes the award, and the Treasury pays it ministerially, from a permanent appropriation, with no further request or approval. It is not a judgment. In 1792 federal judges on circuit declined to decide veterans' pension claims that the Secretary of War could then revise, and some of them wrote to President Washington that the work wasn't the judicial power \autocite{fjc1792hayburn}. Courts decide cases; they don't run programs. A court enters only if a certified award goes unpaid.

The appropriation is permanent because the New Deal's art program died of the other kind. The Federal Art Project paid artists through the Works Progress Administration, and under congressional pressure and budget cuts the WPA, with all its remaining art projects, was dissolved in 1943 \autocite{house2025fap}. An annual appropriation is a vote taken every year on whether the thing should go on, and the vote that ends it needn't say so. A permanent one ends only by a repeal, in the open.
